\newpage
\section{Routing Strategies}
\label{sec:strategies}

% Chapter purpose: Systematic investigation revealing a structural ceiling for
% pre-execution routing on the no-HyDE config space {C0, C2, C3, C6}.
% Narrative: survey (S1–S5, nothing beats C0) → explorations (E1–E3) →
% ceiling analysis (the contribution).
% Data: hotpot/7-routing/ (exploratory), hotpot/8-routing/ (controlled).


\subsection{Experimental Setup}
\label{subsec:routing-setup}

Previously we formulated the routing decision, constructed
oracle label  and assessed the pre-execution signal.  Now, we
evaluate whether learned routing strategies can exploit that signal
beyond a static baseline. The objective remains:
%
\begin{quote}
Given a query~$q$ with pre-execution features
$\mathbf{x} \in \mathbb{R}^d$, select the configuration
$\hat{c}(q) \in \mathcal{C} = \{\text{C0, C2, C3, C6}\}$
that minimises energy while satisfying the quality floor.
\end{quote}
%
This uses no information produced by the pipeline itself. The 
configuration space $\mathcal{C}$ is reduced to non-HyDE configurations.

\begin{figure}[h!]
\centering
\includegraphics[width=0.85\textwidth]{router.png}
\caption{Router flow illustration}
\label{fig:router-flow}
\end{figure}

\paragraph{Evaluation protocol.}
Let $\mathcal{D} = \{(\mathbf{x}_i, y_i)\}_{i=1}^{N}$ denote the routing
dataset, where $\mathbf{x}_i \in \mathbb{R}^{d}$ is the pre-execution
feature vector and $y_i = \ell(q_i) \in \mathcal{C}$ the oracle label
for query~$i$. Features are query-only meaning that every configuration
sees the same~$\mathbf{x}_i$. The label is config-dependent because $\ell(q)$ 
selects the cheapest configuration that satisfies the quality floor:
encoding both energy and quality information in a single target.

The base dataset contains $N = 1{'}200$ queries
with $d = 23$ features; individual experiments vary~$d$ as noted.
All models are evaluated under 5-fold stratified cross-validation,
stratified on~$y_i$ to preserve class proportions,
with 80/20\% train-test splits. Five metrics measure success:

\begin{itemize}
    \item \textbf{Config accuracy}: the fraction of queries routed to the oracle-optimal configuration. Primary metric.
          $$\text{acc} = \frac{1}{N} \sum_{i=1}^{N}
          \mathbb{1}[\hat{y}_i = y_i]$$
    \item \textbf{Config F1 macro}: the macro-averaged F1 across configuration classes.  Penalises
          models that ignore minority classes.
          $$\text{F1}_{\text{macro}} = \frac{1}{|\mathcal{C}|}
          \sum_{c \in \mathcal{C}} \text{F1}_c$$
    \item \textbf{Pass rate}: the fraction of queries whose routed configuration produces an
          acceptable response ($a = 1$).  Pass rate inflates apparent
          performance because multiple configurations satisfy the same query.
          $$\text{pass} = \frac{1}{N} \sum_{i=1}^{N}a(q_i, \hat{y}_i)$$
    \item \textbf{Mean energy/token}: the average energy/token per query under the routed policy.
          $$\bar{E} = \frac{1}{N} \sum_{i=1}^{N} E(q_i, \hat{y}_i)$$
    \item \textbf{Mean ROC AUC} measures a model's average ability to distinguish
    between classes across multiple folds. 
\end{itemize}

\paragraph{Baselines.}
The simplest reference is a \emph{configuration baseline}: apply the same
configuration to every query without any routing.
%$ \hat{c}(q) = c_j,\; \forall q$.
There are four such baselines, one per configuration in~$\mathcal{C}$.
The \emph{oracle} which assigns each query to its most energy-efficient configuration,
$\hat{c}(q) = y_i$. The table below reports the
results.  Among them, C0 achieves the highest accuracy
(0.382) because C0 is the most frequent oracle label.  Any learned router
must exceed this threshold to justify its complexity.
Figure~\ref{fig:baseline-energy} illustrates the energy cost between the baselines.

\begin{table}[h!]
\centering
\begin{tabular}{l r r r r r}
\hline
\textbf{Baseline} & \textbf{Accuracy} & \textbf{F1 macro}  & \textbf{Pass rate} & \textbf{Energy (J)} \\
\hline
Oracle  & 1.000 & 1.000   & 0.555 & 341 \\
C0      & 0.382 & 0.138   & 0.318 & 318 \\
C3      & 0.318 & 0.121   & 0.323 & 324 \\
C2      & 0.153 & 0.066   & 0.363 & 424 \\
C6      & 0.148 & 0.064   & 0.362 & 435 \\
\hline
\end{tabular}
\caption{Configuration baselines and oracle.}
\label{tab:baselines}
\end{table}

\begin{figure}[h!]
\centering
\includegraphics[width=0.7\textwidth]{fig_baseline_energy.pdf}
\caption{Mean energy per query across configuration baselines.}
\label{fig:baseline-energy}
\end{figure}

\paragraph{Investigation design.}
The investigation proceeds in two phases:
\begin{enumerate}
    \item \textbf{Strategy survey}
          (Section~\ref{subsec:strategy-survey})
          Five ML framings of the routing task---multiclass
          classification, binary decomposition, pairwise
          preference, selective abstention and utility scoring---are
          evaluated using the same 23~pre-execution features and
          tree-based classifiers.
    \item \textbf{Explorations}
          (Section~\ref{subsec:explorations})
          We fix the best formulation (utility scoring) and probe
          three dimensions: feature selection~(E1), quality-metric
          definition~(E2) and test-architecture~(E3).
\end{enumerate}

Together, the two phases test whether the performance ceiling is signal-limited or model-limited.

% The controlled study spans four dimensions:
% \begin{itemize}
%       \item 3 feature sets (10~top-importance, 23~handcrafted, 23~+~5~C0 retrieval statistics)
%       \item 2 label definitions (fc~+~faithfulness; fc~+~faithfulness~+~sup\_f1)
%       \item 3 model families (tree-based classifiers, MLP with energy-aware loss (Neural network), two-stage cascade)
%       \item 2 decision architectures (flat two-model with independent hybrid and reranking heads 
%       vs.\ two-stage cascade where Stage~1 predicts search type, executes retrieval and 
%       Stage~2 predicts reranking from the retrieved results).
% \end{itemize}



\subsection{Strategy Survey}
\label{subsec:strategy-survey}

We tested five routing formulations
(S1--S5) to find which prediction target can selectively assign
expensive configurations when they help, rather than defaulting to
the cheapest option for every query. All five use the same 23~pre-execution features,
the same evaluation protocol and tree-based classifiers where
applicable. 
% The table below summarises results whereas 
% the paragraphs below trace the logical progression.

\begin{table}[h!]
\centering
\small
\begin{tabular}{c l l l r}
\hline
\textbf{ID} & \textbf{Strategy} & \textbf{Prediction target} & \textbf{Model} & \textbf{Acc} \\
\hline
S1 & Multiclass          & Oracle label (4-class)        & LR / RF / XGB          & 0.297 \\
S2 & Binary decomp.      & Hybrid? + Rerank? (indep.)    & RF (two heads)         & 0.298 \\
S3 & Pairwise pref.      & Config~A $>$ B? (6 pairs)     & RF + voting            & 0.331 \\
S4 & Selective abstention & Oracle label + confidence     & RF + threshold         & ---   \\
S5 & Utility scoring     & $\hat{P}$(pass $\mid$ q, c)   & RF (per-config) + $\lambda$ & 0.395 \\
\hline
\end{tabular}
\caption{Strategy survey results.}
\label{tab:strategy-survey}
\end{table}

\paragraph{S1: Direct multiclass classification.}
The most direct and naïve approach predicts the configuration label among
\{C0, C2, C3, C6\} using logistic regression, random forest and
XGBoost.  The best classifier (XGBoost) reaches the poor config accuracy of
0.297.

\begin{figure}[h!]
\centering
\includegraphics[width=0.55\textwidth]{fig_cm_s1_multiclass.pdf}
\caption{S1 confusion matrix (XGBoost, best classifier).}
\label{fig:cm-s1}
\end{figure}

\paragraph{S2: Binary decomposition.}
To reduce class imbalance, we decompose the 4-class problem into two
independent binary heads: ``needs hybrid search?'' and ``needs
reranking?''  Two random-forest classifiers predict each decision
separately; their outputs combine into a configuration.  Best accuracy
reaches 0.298. The individual heads are too
weak to improve over majority-class prediction and the selection
mechanism has no notion of energy cost
(Figure~\ref{fig:cm-s2}).

\begin{figure}[h!]
\centering
\includegraphics[width=0.75\textwidth]{fig_cm_s2_binary.pdf}
\caption{S2 confusion matrices for the hybrid and reranking binary heads (RF).}
\label{fig:cm-s2}
\end{figure}

\paragraph{S3: Pairwise preference.}
Instead of predicting a single label, we train six pairwise classifiers
(one per configuration pair) and aggregate via Borda voting.  Best
accuracy reaches 0.331.  When individual pairs cannot reliably
distinguish configurations with 2--3~pp quality differences, voting
amplifies noise rather than signal
(Figure~\ref{fig:cm-s3}).

\begin{figure}[h!]
\centering
\includegraphics[width=0.85\textwidth]{fig_cm_s3_pairwise.pdf}
\caption{S3 pairwise preference confusion matrices (RF).}
\label{fig:cm-s3}
\end{figure}

\paragraph{S4: Selective abstention.}
We add a confidence threshold: route only when the classifier's
predicted probability exceeds a minimum, otherwise fall back to the
cheapest configuration.  At every threshold that maintains meaningful
routing coverage, 68--90\% of queries fall back to the default.
The classifier is rarely confident because the underlying signal is
weak, making abstention impractical
(Figure~\ref{fig:s4-abstention}).

% \begin{figure}[h!]
% \centering
% \includegraphics[width=0.6\textwidth]{fig_s4_abstention_curve.pdf}
% \caption{S4 abstention rate vs config accuracy ($\lambda = 0.75$, fallback C0).}
% \label{fig:s4-abstention}
% \end{figure}

\subsection{Utility Scoring}
\label{subsec:utility-scoring}

The four previous strategies share two failure modes: direct label
prediction is trapped by class imbalance, and no approach accounts
for energy cost during selection.  Utility scoring addresses both.
Instead of predicting the oracle label, we train a binary classifier
\emph{per configuration} that estimates $\hat{P}(\text{pass} \mid q, c)$
---the probability that query~$q$ produces an acceptable response under
configuration~$c$.  Two separate classifiers predict the search-type and
reranking decisions independently.  Configuration selection uses:
%
\begin{equation}
c^*(q) = \arg\max_{c \in \mathcal{C}} \;\;
    \hat{P}(\text{pass} \mid q, c) \;-\; \lambda \cdot
    \widetilde{E}(c)
\label{eq:utility}
\end{equation}
%
where $\widetilde{E}(c)$ is the min-max normalised energy cost of
configuration~$c$ and $\lambda$ controls the energy--quality trade-off.
Ties are broken in favour of the cheaper configuration.  This framework
decouples quality estimation from configuration selection: the classifiers
learn pass probabilities without knowing the energy landscape, while
the utility function injects energy awareness at inference time.

\begin{figure}[h!]
\centering
\includegraphics[width=0.55\textwidth]{fig_cm_s5_utility.pdf}
\caption{S5 confusion matrix (RF, $\lambda = 0.75$).}
\label{fig:cm-s5}
\end{figure}

\paragraph{Survey outcome.}
Utility scoring is the only formulation that avoids the majority-class
trap and introduces energy awareness.  At $\lambda = 1.50$, it exceeds
Static~C0 for the first time, on a minimal note though
(0.395 in the explorations below).  We also tested single-head MLPs
with an energy-aware loss ($\mathcal{L} = \text{BCE} + \alpha \cdot
\delta \cdot \text{sigmoid}(\text{logit})$), but 2{'}000~training
samples proved insufficient: the best variant reached 0.383, matching
the tree baselines without improving on~S5.

Table~\ref{tab:all-strategies} consolidates the best policy from each
strategy alongside the baselines.  All approaches cluster in a narrow
band: config accuracy does not exceed 40\%.

\begin{table}[h!]
\centering
\small
\begin{tabular}{l l r r r r}
\hline
\textbf{Strategy} & \textbf{Best policy} & \textbf{Acc} & \textbf{F1 macro} & \textbf{Pass rate} & \textbf{Energy (J)} \\
\hline
Oracle                 & ---                      & 1.000 & 1.000 & 0.555 & 341 \\
\hline
S5: Utility scoring    & RF $\lambda$=1.50        & 0.395 & 0.189 & 0.328 & 318 \\
MLP (energy loss)      & MLP $\lambda$=1.50       & 0.383 & 0.165 & 0.318 & 317 \\
\hline
Static C0              & ---                      & 0.382 & 0.138 & 0.318 & 318 \\
\hline
S3: Pairwise pref.     & RF + voting              & 0.331 & 0.149 & 0.323 & 335 \\
S2: Binary decomp.     & RF (two heads)           & 0.298 & 0.103 & 0.316 & 361 \\
S1: Multiclass         & XGB                      & 0.297 & 0.131 & 0.325 & 340 \\
\hline
\end{tabular}
\caption{Consolidated strategy results (best policy per approach).}
\label{tab:all-strategies}
\end{table}

\paragraph{Lambda dominance.}
Across all experiments, the best results appear at $\lambda = 1.25$--$1.50$.
Without the energy penalty ($\lambda = 0$), config accuracy drops to
0.28--0.35 as models route queries to expensive configurations.
The energy weight $\lambda$ is more impactful than the choice of
classifier or formulation.  High~$\lambda$ enforces cheap defaults
unless the quality signal is strong, which amounts to a soft version
of Static~C0.  Figure~\ref{fig:lambda-sweep} shows the monotonic
relationship: all three classifiers start 8--10~pp below Static~C0 at
$\lambda = 0$ and converge toward it as $\lambda$ increases.  Only RF
at $\lambda \geq 1.25$ crosses the baseline.

\begin{figure}[h!]
\centering
\includegraphics[width=0.7\textwidth]{fig_lambda_sweep.pdf}
\caption{Config accuracy versus energy penalty $\lambda$}
\label{fig:lambda-sweep}
\end{figure}

\subsection{Routing Explorations}
\label{subsec:explorations}

Given that utility scoring is the most promising formulation that exceeds
Static~C0, we explored some of its boundaries.  Three experiments probe
different dimensions of the approach: which features carry routing
signal~(E1), which quality-metric definition produces learnable
labels~(E2) and whether a two-stage architecture can access
post-retrieval information~(E3).


\subsubsection{Feature Selection (E1)}
\label{subsubsec:feature-importance}

We select the ten features with the highest importance scores from a
preliminary correlation analysis and compare per-model performance
against full 23-feature set.  The goal is to determine whether
the routing signal is concentrated or distributed.

\begin{figure}[h!]
\centering
\includegraphics[width=0.6\textwidth]{fig_e1_feature_importance.pdf}
\caption{Top-10 feature importance scores.}
\label{fig:e1-importance}
\end{figure}

It turns out that the top four are: presence of a logical operator,
embedding similarity gap between the two nearest semantic clusters,
content-word ratio and multi-entity flag.
Structural query properties and embedding-derived difficulty signals
carry the most routing-relevant information for the HotPotQA dataset.

\begin{figure}[h!]
\centering
\includegraphics[width=0.55\textwidth]{fig_e1_feature_comparison.pdf}
\caption{Config accuracy: 10 vs 23 features per model.}
\label{fig:e1-comparison}
\end{figure}

This result is expected but requires empirical verification; the response to feature reduction is model-dependent.
LogReg benefits from noise reduction. RF shows the opposite pattern,
despite ``improving'' from 0.382 to 0.395 with the full
set. XGB is indifferent by default.
The low discriminative power of the
10 features confirms that routing signal
is not limited by noise from the remaining features but
by the features' weak prediction capability. Their mean correlation with the oracle label remains below 0.4.


\subsubsection{Metric Sensitivity (E2)}
\label{subsubsec:metric-selection}

The oracle labels depend on which quality metrics define ``pass'' and
at what threshold.  We train a router (RF, 23~features, $\lambda$~sweep)
under each of the seven non-empty subsets of \{fc, faithfulness,
Sup~F1\}, sweeping $\tau \in \{0.5, 0.6, 0.7, 0.8, 0.9\}$.

\begin{figure}[h!]
\centering
\includegraphics[width=0.85\textwidth]{fig_metric_selection.pdf}
\caption{Config accuracy vs $\tau$ per metric subset}
\label{fig:metric-selection}
\end{figure}

Three regimes emerge in Figure~\ref{fig:metric-selection}:  
\begin{itemize}
      \item At low $\tau$ the oracle is C0-dominated and the router trails its own C0 baseline by 1--2~pp. 
      \item Any combination containing Sup~F1 collapses at $\tau \geq 0.6$ (mean Sup~F1 = 0.545, below threshold).
      \item The fc~+~faith combination occupies a middle regime: at $\tau = 0.7$ the pass rate is 56\% and the
      label distribution is most balanced (C0~38\%, C3~32\%, C2/C6~15\%),
      producing the largest positive margin in the sweep: 0.395 vs
      Static~C0 0.382 (+1.3~pp).  
      \item At stricter thresholds the margin vanishes.
\end{itemize}   

Label learnability depends on the pass-rate regime: too lenient or too
strict both collapse routing to the static baseline.  Noisy quality
metrics (Chapter~\ref{sec:labeling}) further limit the signal any
classifier can exploit (Section~\ref{subsec:ceiling}).


\subsubsection{Cascade Architecture (E3)}
\label{subsubsec:cascade}

All previous strategies commit to the full configuration before
any pipeline component executes.  E3 tests whether deferring the
reranking decision to \emph{midway}---after retrieval---unlocks
additional signal.  Stage~1 predicts search type from query
features (identical to S5); the system executes retrieval; Stage~2
observes five retrieval-quality statistics and decides whether
reranking is needed.  Two variants are tested: V1 uses retrieval
features only (5), V2 adds the 23~query features (28 total).

V1 achieves 0.394 and V2 achieves 0.397---only 0.2~pp above S5.
The Stage~2 reranking classifier reaches AUC~0.53 with retrieval
features alone, barely above random.
Figure~\ref{fig:e7-importance} confirms: three of five retrieval
features rank in the top~10 by RF importance, yet importances are
nearly uniform---a pattern typical of low-signal classification.
The midway signal describes retrieval confidence, not reranking
benefit.

\begin{figure}[h!]
\centering
\includegraphics[width=0.6\textwidth]{fig_e7_feature_importance.pdf}
\caption{E3 Stage~2 feature importance}
\label{fig:e7-importance}
\end{figure}


\subsection{The Structural Ceiling}
\label{subsec:ceiling}

The convergence of all approaches to the 0.38--0.40 range is itself a
finding.  Five strategies and three explorations spanning tree-based
classifiers, neural networks and a two-stage cascade all reach the
same region.
The between-experiment spread is 2~pp (0.377--0.397), comparable to the
within-experiment cross-validation standard deviation of 1--2~pp.
Experiment-to-experiment variation is indistinguishable from fold-to-fold
noise.  A McNemar test on the 1{'}200 query-level predictions confirms
this: the difference between S5~RF (0.395) and Static~C0 (0.382) yields
$p = 0.105$ (two-sided exact test, 86~discordant pairs).  The 1.3~pp
margin is not statistically significant at $\alpha = 0.05$ and translates
to 16~additional correctly-routed queries out of 1{'}199.  This indicates a
structural ceiling, not a failure of any individual model.  E2's metric-selection sweep reinforces this: across
35 label definitions (7 metric subsets $\times$ 5 thresholds), only
fc~+~faith at $\tau = 0.7$ produces a positive margin over Static~C0.
All other combinations either collapse to the baseline or track below
it, indicating that the labels themselves carry noise from imperfect
quality metrics.

\paragraph{Lambda as evidence.}
The $\lambda$ analysis in Section~\ref{subsec:strategy-survey} reinforces
this interpretation.  At $\lambda = 0$ (no energy penalty), all models
achieve 0.28--0.35---worse than Static~C0.  Performance rises only as
$\lambda$ increases toward 1.50, which pushes selections toward cheap
configurations unless the classifier is confident.  In practice, this
amounts to a soft version of ``always pick C0'': the utility function,
not the classifier, drives accuracy above baseline.  The ML model's
marginal contribution is the 1--2~pp it adds on top of the $\lambda$-driven
default.

\paragraph{Weak pre-execution signal.}
The strongest feature--label association is Cramér's $V = 0.152$
(\texttt{q\_has\_logical\_operator}, Section~\ref{subsec:features-signal}).
Cross-validated regression on all 23~features yields
$R^2 \approx 0$.  The features capture query-level properties---length,
structure, named entities---but not the query--document interactions that
determine which pipeline modules help.

\paragraph{Small quality differences.}
Factual correctness varies by less than 0.03~points across
\{C0, C2, C3, C6\}.  When the quality surface is nearly flat, any
classifier faces a low signal-to-noise ratio.  The 95\% confidence interval
within a single configuration covers most of the cross-configuration range.

\paragraph{Module-activation vs.\ model routing.}
The quality gap between configurations (2--4~pp) is an order of magnitude
smaller than the gap between models in model routing (10--30~pp),
fundamentally limiting the routing signal.

\paragraph{The information gap.}
The gap between the oracle (100\%) and the best router (39.5\%) quantifies
the information that exists at retrieval time but is invisible before
execution.  E3 confirmed this directly: even after observing
retrieval-quality statistics, the Stage~2 classifier achieves
AUC $\approx 0.53$ for predicting reranking need---barely above random.

Section~\ref{subsec:ceiling-interpretation} discusses the implications of
these observations.


% \subsection{Chapter Summary}
% \label{subsec:ch7-summary}

% Five routing strategies and four controlled experiments were evaluated
% on 1{'}200~HotpotQA queries using 5-fold cross-validation over the
% four no-HyDE configurations.  Three findings emerge:

% \begin{enumerate}
%     \item All approaches converge to a config accuracy of 0.38--0.40,
%           bounded by a structural ceiling.  No combination of features,
%           quality metrics or model architecture breaks through.
%     \item The highest config accuracy observed (0.397, E4 cascade)
%           exceeds Static~C0 by 1.5~pp---a gain within cross-validation
%           noise.  No approach produces a statistically meaningful
%           improvement over the cheapest-configuration baseline.
%     \item The energy penalty parameter $\lambda$ is more impactful than
%           the choice of classifier, feature set or label definition.
%           High $\lambda$ values enforce cheap defaults unless the quality
%           signal is strong, effectively implementing a soft version of
%           Static~C0.
% \end{enumerate}

% The ceiling is signal-limited, not model-limited.  Module-activation
% routing operates on a nearly flat quality surface where pre-execution
% features cannot capture the query--document interactions that determine
% which pipeline modules help.
% Chapter~\ref{sec:discussion} interprets these findings and discusses
% their implications for future research.
